\subsection{Capteur virtuel du P80 par XGBoost}
\label{sec:xgboost}

\paragraph{Objectif.} Estimer en continu le P80 de la surverse des hydrocyclones, c'est-à-dire la finesse du produit envoyé en aval, sans attendre une analyse granulométrique.

\paragraph{Données.} Le modèle est entraîné sur les données du jumeau (dossier \texttt{data/}) : 2\,880 mesures procédé acquises toutes les 15 minutes du 1\textsuperscript{er} au 30 juillet 2026, et 43\,200 mesures de santé acquises chaque minute. Chaque mesure procédé est associée à la \textbf{moyenne des signaux de santé sur les 15 minutes qui la précèdent}, ce qui reprend le principe d'alignement temporel décrit à la section~\ref{sec:asymetrie}. Sur cette période, le P80 de la surverse varie entre 135,7 et 149,9~µm (moyenne 142,7~µm, écart-type 3,67~µm).

\paragraph{Choix des variables et prévention des fuites.} Un capteur virtuel ne doit utiliser que des grandeurs réellement disponibles au moment de l'estimation. Toutes les granulométries et teneurs mesurées en aval, ainsi que le rapport de réduction (calculé à partir du P80), ont donc été \textbf{exclues}, car elles contiennent la réponse. Deux scénarios ont été évalués :
\begin{itemize}
  \item \textbf{Scénario A -- capteurs en ligne} (32 variables) : débits et fractions solides des différentes lignes (débitmètres et densimètres), ainsi que les signaux de santé des équipements (puissance et vitesse du broyeur, pressions des cyclones, vitesse de la pompe, niveau du bac, etc.) ;
  \item \textbf{Scénario B -- A + caractérisation de l'alimentation} (34 variables) : on ajoute le P80 et la teneur BPL de l'alimentation fraîche, connus en amont du circuit.
\end{itemize}

\paragraph{Protocole.} Pour respecter la nature temporelle des données, le découpage est \textbf{chronologique} : 70\,\% pour l'entraînement (2\,015 points), 15\,\% pour la validation (433 points) et 15\,\% pour le test (432 points, du 26 au 30 juillet). Un découpage aléatoire aurait laissé le modèle «~voir le futur~» et surestimé ses performances. Les hyperparamètres de XGBoost~\cite{chen2016} (profondeur, taux d'apprentissage, sous-échantillonnage) sont choisis par recherche de grille, avec arrêt anticipé sur le jeu de validation. Le modèle final est ensuite réentraîné sur l'entraînement et la validation, puis évalué une seule fois sur le test. Trois modèles servent de comparaison :
\begin{itemize}
  \item une prédiction constante égale à la moyenne, qui sert de plancher ;
  \item une régression linéaire régularisée (Ridge) ;
  \item un Random Forest avec les mêmes hyperparamètres que le simulateur What-If du backend.
\end{itemize}
Les performances sont mesurées par le coefficient de détermination $R^2$, la \gls{rmse}, la \gls{mae} et l'erreur relative moyenne (MAPE) :
\[
\mathrm{RMSE}=\sqrt{\frac{1}{n}\sum_{i=1}^{n}(y_i-\hat{y}_i)^2}, \qquad
\mathrm{MAE}=\frac{1}{n}\sum_{i=1}^{n}\lvert y_i-\hat{y}_i\rvert, \qquad
R^2 = 1-\frac{\sum_i (y_i-\hat{y}_i)^2}{\sum_i (y_i-\bar{y})^2}
\]

\paragraph{Résultats.} Le tableau~\ref{tab:resp80} et la figure~\ref{fig:r2comp} présentent les performances sur le jeu de test.

\begin{table}[H]
\centering
\caption{Performance des modèles sur le jeu de test (432 points, 26--30 juillet 2026)}
\label{tab:resp80}
\small\setlength{\tabcolsep}{5pt}
\begin{tabular}{llcccc}
\toprule
\textbf{Scénario} & \textbf{Modèle} & $\boldsymbol{R^2}$ & \textbf{RMSE (µm)} & \textbf{MAE (µm)} & \textbf{MAPE} \\
\midrule
\multirow{4}{*}{A -- en ligne}
 & Moyenne (référence) & $-0{,}025$ & 3,916 & 3,542 & 2,50\,\% \\
 & Régression Ridge & 0,890 & 1,281 & 1,002 & 0,70\,\% \\
 & Random Forest & 0,953 & 0,842 & 0,672 & 0,47\,\% \\
 & \textbf{XGBoost} & \textbf{0,950} & \textbf{0,869} & \textbf{0,669} & \textbf{0,47\,\%} \\
\midrule
\multirow{4}{*}{B -- + alimentation}
 & Moyenne (référence) & $-0{,}025$ & 3,916 & 3,542 & 2,50\,\% \\
 & Régression Ridge & 0,950 & 0,861 & 0,662 & 0,47\,\% \\
 & Random Forest & 0,952 & 0,845 & 0,667 & 0,47\,\% \\
 & \textbf{XGBoost} & \textbf{0,953} & \textbf{0,841} & \textbf{0,657} & \textbf{0,46\,\%} \\
\bottomrule
\end{tabular}
\end{table}

\begin{figure}[H]
\centering
\includegraphics[width=0.85\textwidth]{models_r2_comparison}
\caption{Coefficient de détermination $R^2$ des quatre modèles sur le jeu de test, pour les deux scénarios}
\label{fig:r2comp}
\end{figure}

La figure~\ref{fig:p80ts} superpose le P80 mesuré et le P80 estimé par XGBoost sur toute la période de test, et la figure~\ref{fig:parity} compare les valeurs point par point.

\begin{figure}[H]
\centering
\includegraphics[width=\textwidth]{p80_timeseries_B}
\caption{P80 mesuré et P80 estimé par XGBoost sur le jeu de test (scénario B)}
\label{fig:p80ts}
\end{figure}

\begin{figure}[H]
\centering
\includegraphics[width=0.45\textwidth]{p80_parity_A}\hfill
\includegraphics[width=0.45\textwidth]{p80_parity_B}
\caption{P80 estimé en fonction du P80 mesuré (la diagonale correspond à une estimation parfaite)}
\label{fig:parity}
\end{figure}

\paragraph{Robustesse.} Un seul découpage peut être favorable par hasard. Une validation croisée temporelle en 5 plis (\textit{TimeSeriesSplit}), où chaque pli teste le modèle sur une période postérieure à celle d'entraînement, donne pour XGBoost :
\begin{itemize}
  \item scénario A : $R^2 = 0{,}938 \pm 0{,}008$ et RMSE $= 0{,}93 \pm 0{,}06$~µm ;
  \item scénario B : $R^2 = 0{,}944 \pm 0{,}007$ et RMSE $= 0{,}88 \pm 0{,}03$~µm.
\end{itemize}
La faible dispersion d'un pli à l'autre montre que le modèle reste stable dans le temps.

\paragraph{Variables les plus influentes.} La figure~\ref{fig:importance} montre que, dans le scénario B, le modèle s'appuie d'abord sur la caractérisation de l'alimentation : P80 d'alimentation (55\,\% du gain total), teneur BPL (21\,\%) et débit solide d'alimentation (13\,\%). Dans le scénario A, en l'absence de ces analyses, il reconstitue l'information à partir des mesures hydrauliques : fraction solide de la surverse (36\,\%), débit de décharge du broyeur (23\,\%), débit de surverse (16\,\%) et débit d'alimentation (15\,\%). Ce comportement est cohérent avec la physique de la classification : la densité et le débit de la surverse dépendent directement de la coupure granulométrique des hydrocyclones~\cite{wills2016}.

\begin{figure}[H]
\centering
\includegraphics[width=0.75\textwidth]{xgb_feature_importance}
\caption{Importance des variables du modèle XGBoost (scénario B, gain normalisé)}
\label{fig:importance}
\end{figure}

\paragraph{Analyse.} Quatre enseignements se dégagent.
\begin{enumerate}
  \item \textbf{Le capteur virtuel est fiable.} Avec une erreur moyenne de 0,66~µm, soit 0,46\,\% du P80, XGBoost réduit l'erreur de 81\,\% par rapport à la référence constante (MAE de 3,54~µm).
  \item \textbf{Les capteurs en ligne suffisent.} Sans aucune analyse de laboratoire (scénario A), XGBoost atteint déjà $R^2 = 0{,}950$. Un modèle linéaire plafonne à 0,890 dans ce scénario : la relation entre les débits et la granulométrie n'est pas linéaire, ce qui justifie un modèle à base d'arbres. À l'inverse, dans le scénario B, la régression Ridge atteint 0,950, presque autant que XGBoost (0,953) : lorsque la caractérisation de l'alimentation est connue, un modèle linéaire suffit. L'intérêt de XGBoost se situe donc dans le cas le plus réaliste, celui où l'on ne dispose que des capteurs en ligne.
  \item \textbf{XGBoost et Random Forest sont à égalité en précision.} Les écarts, de l'ordre de 0,003 en $R^2$, ne sont pas significatifs. XGBoost se distingue par sa \textbf{sobriété} : son modèle final comporte 49 arbres de profondeur 3 et pèse \textbf{0,06~Mo}, contre 100 arbres de profondeur 22 et \textbf{8,3~Mo} pour le Random Forest, soit un modèle environ 130 fois plus léger. C'est un atout pour un déploiement embarqué, au plus près des machines.
  \item \textbf{Le $R^2$ reflète surtout la reconnaissance du type de minerai.} Les figures~\ref{fig:p80ts} et~\ref{fig:parity} montrent que le P80 évolue selon trois régimes (environ 138, 141,5 et 147~µm), qui changent en moyenne toutes les 3,3~heures et suivent la nature du minerai alimenté (P80 et teneur BPL de l'alimentation). Le modèle identifie chaque changement de régime sans retard. En revanche, \textbf{à l'intérieur d'un régime}, il n'explique pas les fluctuations : le $R^2$ calculé sur les écarts à la moyenne de chaque régime est proche de zéro ($-0{,}16$). Ces fluctuations ne sont donc liées à aucune variable d'entrée et se comportent comme un bruit de mesure, ce qui explique aussi que les trois modèles non linéaires convergent vers la même RMSE, environ 0,84~µm. Il serait donc trompeur de présenter le $R^2$ de 0,95 comme une capacité à suivre chaque micro-variation : sa valeur opérationnelle tient à la détection immédiate des changements de minerai, évaluée ci-dessous.
\end{enumerate}

\paragraph{Reproductibilité.} L'ensemble de l'expérience tient dans un script unique du projet, \nolinkurl{ml/soft_sensor_p80/train_xgboost.py}, avec une graine aléatoire fixée. Le script produit le fichier \texttt{metrics.json}, les figures de cette section et le modèle entraîné (\nolinkurl{xgb_soft_sensor_p80.joblib}).

\paragraph{Comparaison avec la pratique actuelle : l'analyse de laboratoire.} Un $R^2$ ne dit pas si l'outil est utile à l'exploitant. La vraie question est : \textit{le capteur virtuel fait-il mieux que ce dont l'opérateur dispose aujourd'hui ?} Sans granulomètre en ligne, le P80 est connu par une analyse de laboratoire périodique. Nous avons simulé cette pratique sur la même période de test : un prélèvement toutes les 4 ou 8 heures, un résultat disponible une heure après le prélèvement, et, entre deux résultats, l'opérateur se fie à la dernière valeur connue. Le capteur virtuel est ici le modèle XGBoost du scénario A, qui n'utilise que les capteurs en ligne (tableau~\ref{tab:labo}, figures~\ref{fig:labo} et~\ref{fig:labomae}).

\begin{table}[H]
\centering
\caption{Capteur virtuel et analyse de laboratoire sur le jeu de test (432 points, 26--30 juillet 2026)}
\label{tab:labo}
\small
\begin{tabularx}{\textwidth}{>{\raggedright\arraybackslash}X ccc}
\toprule
 & \textbf{Labo / 8 h} & \textbf{Labo / 4 h} & \textbf{Capteur virtuel} \\
\midrule
Erreur absolue moyenne & 4,26~µm & 3,26~µm & \textbf{0,66~µm} \\
Erreur maximale & 11,2~µm & 10,9~µm & \textbf{4,4~µm} \\
Part du temps avec une erreur $> 2$~µm & 61\,\% & 50\,\% & \textbf{2\,\%} \\
Régime de minerai correctement identifié & 39\,\% & 52\,\% & \textbf{98,6\,\%} \\
\bottomrule
\end{tabularx}
\end{table}

\begin{figure}[H]
\centering
\includegraphics[width=\textwidth]{p80_capteur_vs_labo}
\caption{48 premières heures du test : P80 réel, capteur virtuel avec son intervalle à 90\,\%, et valeur connue de l'opérateur avec une analyse de laboratoire toutes les 8~h}
\label{fig:labo}
\end{figure}

\begin{figure}[H]
\centering
\includegraphics[width=0.62\textwidth]{p80_mae_labo_vs_capteur}
\caption{Erreur absolue moyenne sur le P80 selon la source d'information}
\label{fig:labomae}
\end{figure}

Le résultat est net : l'erreur est \textbf{divisée par 6,5} par rapport à une analyse toutes les 8~heures, et par 5 par rapport à une analyse toutes les 4~heures. La raison est visible sur la figure~\ref{fig:labo} : le minerai change de nature toutes les 3~heures environ, plus vite que le laboratoire ne peut le suivre. L'opérateur pilote alors une bonne partie du temps avec une valeur périmée, souvent celle d'un autre régime. Le capteur virtuel, lui, se recale à chaque mesure, soit toutes les 15~minutes. Sur les 432 points du test, il ne se trompe de régime que 6 fois, toujours au moment même d'une transition.

\paragraph{Intervalle de confiance de chaque estimation.} Un opérateur doit savoir quand faire confiance à une estimation. Nous avons donc associé à chaque valeur un \textbf{intervalle de prédiction à 90\,\%}, construit par prédiction conforme (\textit{split conformal prediction})~\cite{angelopoulos2023}. Le principe est simple : on mesure les erreurs du modèle sur le jeu de validation, qu'il n'a pas vu pendant l'apprentissage, et l'on retient l'erreur que 90\,\% d'entre elles ne dépassent pas. On obtient un intervalle de \textbf{$\pm$\,1,41~µm}. Sur le jeu de test, \textbf{91,4\,\%} des valeurs réelles tombent dans cet intervalle, conformément à la garantie annoncée de 90\,\%. Le capteur virtuel peut ainsi afficher, par exemple, «~P80 = 147,2~µm $\pm$ 1,4~µm~».

\paragraph{Reproductibilité.} Ces analyses sont produites par le script \nolinkurl{ml/soft_sensor_p80/analyse_avancee.py}.

